\documentclass[10pt,a4paper]{amsart}
\usepackage[margin=19mm]{geometry}
\usepackage{amsmath,amssymb,mathtools}
\usepackage[scr=rsfs,cal=euler]{mathalfa}
\usepackage{xfrac}
\usepackage[unicode,hidelinks,bookmarks=false]{hyperref}
\newcommand{\ie}{i.e.\ }
\newcommand{\cf}{cf.\ }
\newcommand{\resp}{respectively\ }
\newcommand{\Subsection}{\subsection}
\providecommand{\nobreakdash}{\nobreak\hbox{-}}
\setlength{\emergencystretch}{2em}
\multlinegap0pt
\usepackage{bm}
\usepackage{bbm}
\usepackage{dsfont}
\usepackage{xspace}
\usepackage{cancel}
\usepackage{mathtools}
\usepackage{amsthm}
\makeatletter
\@ifundefined{thm}{\newtheorem{thm}{Thm}}{}
\@ifundefined{lem}{\newtheorem{lem}[thm]{Lemma}}{}
\makeatother
\makeatletter
\newcommand{\Sph}{\mathbb{S}}
\expandafter\ifx\csname diagram\endcsname\relax
\newenvironment{diagram}[1][]{%
    \begin{equation}%
    \begin{tikzcd}[#1]%
    \refstepcounter{dia}
}{%
    \end{tikzcd}%
    \end{equation}%
}
\fi
\makeatother
\title[Comparing Periodic Point Invariants for Parameterized Famili]{Comparing Periodic Point Invariants for Parameterized Families of Maps}
\author{Lucas Williams}
\date{Source: arXiv:2508.18339v1 (CC BY 4.0)}
\begin{document}
\maketitle

\noindent\emph{Source and licence.} Comparing Periodic Point Invariants for Parameterized Families of Maps. Version arXiv:2508.18339v1 (CC BY 4.0), distributed under the Creative Commons Attribution 4.0 International licence, as recorded in the arXiv licence metadata for this exact version and shown on the abstract page https://arxiv.org/abs/2508.18339v1. Full rights notice: licenses/arxiv-2508.18339-CC-BY-4.0.txt.\par\smallskip

\noindent\textit{Editorial scope.} Counted unit: the Lem whose source label is \texttt{lem:split section} in the article (source0.tex lines 698--700, source label \texttt{lem:split section}), delivered together with the article's own complete original proof of it (source0.tex lines 702--725, one \texttt{proof} environment of 24 lines). No mathematics is written, repaired or re-derived: statement and proof are the article's own text line for line. Required context retained uncounted: none. Every definition, display and label of those lines is retained unchanged. Omitted from this package and not redistributed: 1--697, 701--701, 726--1101. Nothing inside a retained excerpt is omitted or altered: every retained body is delivered exactly as the article's lines are written, so no omission receipt of any kind accompanies this package. Citations: the retained excerpts carry no citation command at all, so no bibliography entry of the article is transcribed and no reference is redistributed. Reference adaptation: none was needed and none was made; every reference inside the delivered unit resolves inside it, so no unresolved reference or citation is produced. Printed slips kept literal: no printed word, symbol, hypothesis or number of the counted statement or of its proof was altered. Layout and class: the article's own class is \texttt{amsart} with packages including amsmath, amssymb, amsthm, array, color, multirow, pdflscape, mathrsfs, subcaption, todonotes, tikz, tikz-cd, float, imakeidx; the wrapper is the lane's pinned \texttt{amsart} editorial wrapper at 10pt on A4, which restates the theorem-like environments the retained text needs and adds the article's own macro definitions that the retained text needs (\texttt{\textbackslash Sph}), with extra wrapper packages (bm, bbm, dsfont, xspace, cancel, mathtools). The delivered numbering is the unit's own. Assets: no retained span contains a figure environment, a graphics-inclusion command, a PDF-page inclusion, a table or a TikZ picture; no image file is copied into this package, the package declares no asset dependency and no image file is redistributed with it. Licence: version arXiv:2508.18339v1 of this article is distributed under the Creative Commons Attribution 4.0 International licence, as recorded in the arXiv licence metadata for this exact version and shown on the abstract page https://arxiv.org/abs/2508.18339v1. A full rights notice is delivered at licenses/arxiv-2508.18339-CC-BY-4.0.txt. Characters: every retained excerpt is pure ASCII, so the delivered engine, XeLaTeX with its default Latin Modern text font, reports no missing character.
\begin{lem}\label{lem:split section}
The spectrum of sections, $\Gamma_{S^1}\Sigma_{+S^1}^\infty S^1$, is stably equivalent to the cotensor spectrum $F(S^1_+,\Sph)$ and thus, to $\Sph\vee  \Sph^{-1}$. Moreover, the $\Sph$ summand corresponds to sections of $\Sigma_{+S^1}^\infty S^1$ which are constant in each fiber.
\end{lem}

\begin{proof}
Spectrum level $k$ of $\Gamma_{S^1}\Sigma_{+S^1}^\infty S^1$ is the space of sections $\Gamma_{S^1}S^k\overline{\wedge}_{+S^1}S^1$. Moreover, $S^k\overline{\wedge}_{+S^1}S^1 \simeq S^k\times S^1$. Therefore, a section of this space over $S^1$ is just a map $S^1_+\to S^k$. This identification commutes with the structure maps of $\Gamma_{S^1} \Sigma_{+S^1}^\infty S^1$ so that it induces an equivalence 
\[
\Gamma_{S^1}\Sigma_{+S^1}^\infty S^1 \simeq F(S^1_+,\Sph).
\]

Now observe that we have a retract of spaces 
\[
*_+ \to S^1_+ \to *_+.
\]
This induces a retract of cotensor spectra
\[
F(*_+,\Sph) \to F(S^1_+,\Sph) \xrightarrow{j} F(*_+,\Sph).
\]
Therefore, 
\[
F(S^1_+,\Sph)\simeq F(*_+,\Sph)\vee Fj
\]
where $Fj$ is the fiber of $j$. We may identify $F_j$ with $F(S^1,\Sph)$. Then identifying $F(*_+,\Sph)\simeq \Sph$ and $F(S^1,\Sph)\simeq \Sph^{-1}$, we obtain
\[
\Gamma_{S^1} \Sigma_{+S^1}^\infty S^1\simeq \Sph\vee \Sph^{-1}.
\]
The last part of the lemma statement follows from the fact that the map $j$ is induced by mapping $S^1$ to a single point. 
\end{proof}

\end{document}
